\documentclass[11pt]{article}
\usepackage{fiche}
\matiere{Méthodes de scoring}
\chapitre{M2}
\source{Méthodes de scoring.pdf}
%% ocr: cer=0.000; exemples=2; lecture=complex
%% statut: propre

\begin{document}
%% page: Méthodes de scoring.pdf
\titre{\entoure{9} \souligne{Méthodes de scoring}}

\titre{\entoure{a} \souligne{Modélisat° et choix du seuil}}

$\left\{ \begin{array}{l} Y = \text{critère à modéliser} \\ X = \text{caractéristique individuelle} \end{array} \right.$
\entoure{$Y = \begin{cases} 1, & \Pr[Y=1] = p_1 = 1 - p_0 \\ 0, & \Pr[Y=0] = p_0 = 1 - p_1 \end{cases}$}

\begin{attention}
\couleur{rouge}{$p_1 \neq p_0$ = inconnue mais estimable}
\end{attention}

\begin{itemize}
\item \couleur{bleu}{On considère des variables $X$ continue et de densité $f(X)$. \souligne{Les densités conditionnelle $f_1(X)$ et $f_0(X)$ décrivent les caractéristiques individuelle selon l'état de $Y$}}
\end{itemize}

\begin{acote}[10,23,37,30]{deux courbes de densité en cloche qui se chevauchent : $f_{Y=0}(X)$ en vert à gauche, $f_{Y=1}(X)$ en rouge à droite}
\couleur{rouge}{$\left\{ \begin{array}{l} \text{\entoure{+} c'est différente dans f's sont éloignés} \\ \text{\entoure{+} ça induit que les variables ont une caractéristique différente selon l'état de     } Y \Rightarrow \text{    \souligne{meilleur discriminat°}} \end{array} \right.$}
\couleur{bleu}{Donc si c'est confondu $\Rightarrow$ y'a pas de différence qui permettent de discriminer les indiv}
\end{acote}

\begin{itemize}
\item \couleur{bleu}{Par les formules de Bayes, les \souligne{proba conditionnelle} sont :}
\begin{encadre}
\[ \entoure{$\begin{aligned} p_1(X) &= \Pr[Y=1 \mid X=x] \\ p_0(X) &= \Pr[Y=0 \mid X=x] \end{aligned}$} \Leftrightarrow \entoure{$\begin{aligned} p_1(X) &= \frac{p_1 f_1(x)}{f(x)} \\ p_0(X) &= \frac{p_0 f_0(x)}{f(x)} \end{aligned}$} \]
\[ \text{avec } f(x) = p_1 f_1(x) + p_0 f_0(x) \]
\end{encadre}

\item \couleur{bleu}{Imaginons qu'on ne connaît pas les caractéristique individuelles mais qu'on souhaite déterminer la région $\mathcal{A}$ et le seuil du score (le cutoff)}

\couleur{bleu}{$\Rightarrow$ \souligne{Modélisation va essayer de minimiser les erreurs d'affectat°}}
\end{itemize}

\couleur{rouge}{\# Forcément, dans le business, on max la \souligne{rentabilité} \\ Voici les paramètres de rentabilité}

\couleur{bleu}{\entoure{Le gain issu d'une bonne décision accordé à un bon}}
\[ g = \Pr[\underline{\mathbf{1}_{\mathcal{A}}(X)} ; \underline{Y=1}] \]
\couleur{vert}{le score prédit que l'indiv est bon \quad et qu'il l'est vraiment}

\couleur{violet}{\entoure{Le coût d'octroi à un mauvais client}}
\[ c_0 = \Pr[\mathbf{1}_{\mathcal{A}}(X) ; \underline{Y=0}] \]
\couleur{vert}{score prédit que c'est bon \quad alors qu'il ne l'est pas}

\couleur{violet}{\entoure{Le manque à gagner de non-octroi d'un bon client}}
\[ c_1 = \Pr[\mathbf{1}_{\bar{\mathcal{A}}}(X) ; \underline{Y=1}] \]
\couleur{vert}{score ne l'accorde pas \quad alors que c'est un bon client}

\couleur{rouge}{$\hookrightarrow$ Banque va donc maximiser ses gains/rentabilité :}
\couleur{rouge}{\[ \mathcal{A}^* = \underset{\mathcal{A}}{\arg\max}\, R(\mathcal{A}), \quad \text{avec} \]}
\couleur{rouge}{\begin{multline*} R = g \cdot \Pr[\mathbf{1}_{\mathcal{A}}(X) \mid Y=1] \cdot \Pr(Y=1) \\ - c_0 \cdot \Pr(\mathbf{1}_{\mathcal{A}}(X) \mid Y=0) \cdot \Pr(Y=0) \\ - c_1 \cdot \Pr(\mathbf{1}_{\bar{\mathcal{A}}}(X) \mid Y=1) \cdot \Pr(Y=1) \end{multline*}}
\end{document}
